\scanpage{92}
drais l'interpréter comme l'extension \uline{universelle} (de Serre) de $B$ par un fibré vectoriel.

Notons d'abord que le schéma $B^*$ des translations de $\deleted{\unsure{B}}\inserted{A}$ opère \unsure{trivialt} sur l'extension $\deleted{\unsure{E_{A/S}}}\inserted{M^0_{A/S}}$ (car il opère \unsure{trivialt} sur $\underline{\mathrm{Pic}}^0_{A/S}$ et sur $\underline{\omega}_{A/S}$ \deleted{\unsure{\}}}, et $\forall$ hom de $B$ dans $\underline{\omega}_{\unsure{X|A}/S}$ est trivial). Donc par théorie de descente on voit que $M^0_{A/S}$ est can isom à $M^0_{B^*/S}$, donc c'est un invariant qui ne dépend que du \deleted{\illegible{1}} schéma abélien $B$ sur $S$.

Pour calculer la classe de l'extension obtenue de $\underline{B}$ par $\underline{\omega}_{A/S}$ comme élément de
\[
\mathrm{H}^0\bigl(S,\underbrace{\mathrm{R}^1f_{B*}\deleted{\unsure{*}}(\underline{\omega}_{A/S}\otimes_{\mathcal{O}_S}\mathcal{O}_B)}_{\inserted{\unsure{\mathrm{R}^1f_{B*}(\mathcal{O}_B)^\vee}\otimes\mathrm{R}^1f_{B*}(\mathcal{O}_B)}}\bigr)
\]
on regarde \inserted{par la recette générale} le faisceau inversible canonique sur $P=A\times_S B$, et on écrit l'obstruction à mettre dessus une connexion relative à $B$, i.e. la classe de cohomologie qui lui correspond dans $\mathrm{H}^1(P,\underline{\Omega}^1_{P/B})$ [\unsure{pens. à localiser} sur $S$].
\margin{\diagram{\begin{tikzcd}[column sep=small,row sep=small]
A \arrow[dr,dash,"f_A"'] && B \arrow[dl,dash,"f_B"] \\
& S &
\end{tikzcd}}}
\margin{\diagram{\begin{tikzcd}[column sep=0pt,row sep=small]
A & {}\times_S B \\
B^* \arrow[u,equal] &
\end{tikzcd}}}

Calculons donc la classe de cohomologie de $\underline{L}$ dans $\mathrm{H}^1(A\times_S B,\underline{\Omega}^1_{A\times_S B/S})\deleted{\unsure{\simeq\mathrm{H}}}$, localisé sur $S$.
\margin{91}
